\chapter{Calculos justificativos}

\section{Parametros de diseno}
Los parametros base del diseno se resumen a continuacion. Los valores marcados
como supuesto se asumen para el anteproyecto y deben confirmarse con la
concesionaria antes de la construccion.

\begin{center}
\small
\begin{tabularx}{\linewidth}{@{}L l X@{}}
\toprule
Parametro & Valor & Nota \\
\midrule
Tension nominal & 380/220~V, 3F+N+PE & DEC-004 \\
Frecuencia & 60~Hz & sistema peruano \\
Sistema de puesta a tierra & TN-S & DEC-005 \\
Capacidad de servicio & \ServCap & propuesta \\
Factor de reserva & 1,20 & margen de diseno \\
Factor de potencia de demanda & 0,83 & MD 15,67~kW / 18,83~kVA \\
Desbalance de fases & \Desb & objetivo \(\leq\) 10~\% \\
Resistividad del cobre (en servicio) & 0,021~\(\Omega\cdot\)mm\textsuperscript{2}/m & Tabla del calculo \\
Reactancia del conductor & 0,08~\(\Omega\)/km & hipotesis \\
Factor de desrating & 0,80 & CNE-U Tabla 5A/5C: temperatura, agrupamiento y altitud \\
Caida de tension (alim./derivado) & 2,5~\% / 4~\% total & CNE-U Regla 050-102 \\
Conductor minimo & 2,5~mm\textsuperscript{2} & CNE-U Regla 030-002 \\
Icc en punto de entrega & \IccAsumido~kA (supuesto) & pendiente factibilidad \\
Poder de corte del ITM principal & \(\geq\) \Icu~kA & cubre el Icc supuesto \\
Altitud de sitio & 3830~m s.n.m. & confirmada por el estudiante \\
\bottomrule
\end{tabularx}
\end{center}

\section{Demanda y suministro}
La potencia instalada es \PI. La maxima demanda calculada es
\MD~(\MDkw~kW y \MDkva~kVA), con un desbalance de fases de \Desb, por debajo
del objetivo de 10~\%.

Aplicando el factor de reserva de 20~\% la demanda con reserva es \MDReserva,
menor que la capacidad de servicio propuesta de \ServCap. La corriente maxima
de fase con reserva es \Ipase, cubierta por el interruptor principal de
\IPrincipal.

La demanda se determina con cargas realmente proyectadas y factores de demanda
y simultaneidad justificados, conforme al articulo 7 de la EM.010. Para usos
no residenciales se verifica ademas la cota de la Regla 050-210 del CNE-U con
la Tabla 14 (25~W/m\textsuperscript{2} para uso industrial/comercial, factor
de demanda 100~\%), que queda ampliamente superada por las cargas instaladas.

\section{Cortocircuito y selectividad}
Para el anteproyecto se asume una corriente de cortocircuito de
\IccAsumido~kA en el punto de entrega, valor tipico de un suministro BT de
esta magnitud y que debe confirmarse con la factibilidad de la concesionaria.
El interruptor principal de \IPrincipal~se selecciona con poder de corte
\(I_{cu} \geq\) \Icu~kA, por lo que cubre holgadamente el Icc supuesto.

La selectividad entre el interruptor general y los interruptores de los
subalimentadores (TDE \ITMAlTDE~A, TDF \ITMAlTDF~A y TD-A1 \ITMAlTDA~A) se logra
por escalonamiento de corrientes nominales, con relaciones de 1,25 a 2,5~veces
entre niveles, garantizando que solo opere el dispositivo mas cercano al
defecto. El estudio de selectividad definitivo y el calculo de Icc por fase se
realizan en la etapa de construccion con los datos reales de la concesionaria.

\section{Conductores y caida de tension}
La comprobacion usa las capacidades de la Tabla 2 del CNE-U para cobre
XLPE/EPR a 90~\({}^\circ\)C (tres conductores cargados, metodos B1 y D de la
NTP 370.301) y secciones de enlace equipotencial de la Tabla 16. El conductor
minimo es 2,5~mm\textsuperscript{2}, conforme a la Regla 030-002.

Conforme a la Regla 050-102, la caida de tension de cada alimentador o
circuito derivado no supera 2,5~\% y la caida total hasta la salida mas
alejada no supera 4~\%. El alimentador principal y los subalimentadores se
calculan con la formula trifasica (380~V). El alimentador principal presenta
una caida de \dVAlim.

El desrating de 0,80 adopta conservadoramente la correccion por temperatura,
agrupamiento y altitud (Puno, aproximadamente 3830~m), sin aplicar factores
separados.

\section{Cuadro de circuitos}
\ALIMENTADORESTAB

\bigskip
\CIRCUITOSTAB

\section{Coordinacion de arranque}
El arranque de los motores se verifica contra la coordinacion del
subalimentador correspondiente. El arranque del mayor motor de cada tablero
produce corrientes momentaneas que los interruptores de curva D toleran
durante el transitorio; el grupo electrogeno se dimensiona para el arranque
secuencial de una bomba sumergible a la vez. Los valores de catalogo deben
sustituirse por placas antes de construir.

\section{Grupo electrogeno y UPS}
La carga critica permanente es \EmergKVA. El escenario de arranque secuencial
con el mayor motor es \ArranqueKVA y con margen de 1,25 es \ArranqueMargen. A
la altitud de sitio (factor \AltFactor) el grupo \GrupoModelo~de~\GrupoNom~kVA
standby entrega \GrupoDisp, por lo que la verificacion CUMPLE.

Las UPS se dimensionan por sus cargas conectadas: la UPS de
combustible/control y la UPS de POS-CCTV-comunicaciones quedan por encima de
sus potencias nominales de forma holgada.
